% Gerado por tools/tei2latex.py a partir de entrada.xml. Regenerar, não editar.

%%>>> parte-2

%%>>> cap-7

%%>>> art-44

\section*{Artigo I.\\ Creação, sua natureza, principio e fim.}
\addcontentsline{toc}{section}{Artigo I. --- Creação, sua natureza, principio e fim.}

\pg{513} 85. \textbf{Creação}. --- \emph{Creação} é o acto, pelo qual \textsc{Deus} tirou do \emph{nada} todas as coisas. --- Tirar do \emph{nada} o ente é produzil-o na sua \emph{totalidade}, isto é, não só emquanto á \emph{forma}, mas tambem emquanto á \emph{materia} \footnote{Tratamos, neste capitulo, dos principaes actos do amor \emph{efectivo} de Dzvus para com as creaturas, que são: a \emph{creacdo}, 8 conservação, o concurso e a \emph{providencia}. --- Neste artigo occupamo-nos da \emph{creacdo}. (\emph{Cf}. \emph{Sum}. \emph{Tos} p.d; q. 44, 45; \emph{c}. \emph{Gent}, 1. II, c. 15 e seg).\par A \emph{creagüo} é a \emph{produccdo de uma coisa do nada}. Para que, pois, a produecüo de uma coisa possa chamar-se \emph{creação}, são necessarios e sufficientes dois requisitos: 1.º) que \emph{nada preewista} da coisa, que se diz \emph{creada}, a não ser a não-repugnancia para existir, ou a absoluta possibilidade, ou a potencia objectiva, e por isso a creação differe de todas as outras especies de producção, como são a geração, a alteração, etc.; --- 2.9) que a coisa, que se diz creada, passe do não-ser para o ser, não emquanto ella é alguma coisa no termo a \emph{quo}, que é o \emph{não-ser}, mas emquanto, pelo influxo de uma causa proporcionada, recebe o ser, que antes não tinha.\par A \emph{creacüo} diz-se \emph{activa} ou \emph{passiva}, conforme se considera em \textsc{DEus} ou na creatura. --- Considerada em \textsc{Deus}, a \emph{creação} é o acto, pelo qual \textsc{Deus} tira do \emph{nada} as coisas. E' a definição apresentada no texto. Este acto de \textsc{Deus} é a sua propria essencia, emquanto se refere ás creaturas. Para a comprehensüo d'esta verdade, é necessario distinguir um duplice acto --- um \emph{immanente} e outro \emph{transeunte}. \emph{O immanente} é o que permanece no agente, e é uma perfeição do proprio agente, como entender, querer, etc.; o \emph{transeunte} é o que passa para o effeito externo, \emph{e} é uma perfeição do proprio effeito, que existe por elle, como aquecer, fabricar, etc. Este duplice acto compete a \textsc{Deus}. Compete-lhe o primeiro acto, porque \textsc{Deus} entende, quer, -ama, gosa; compete-Ihe o segundo, emquanto produz as coisas, conserva-as e\par 88 VOL. Il.\pgn{514}rege-as. E porque o primeiro acto é perfeição do agente, e o segundo é perfeição do effeito, e o agente é anterior ao effeito e é causa d'elle, é evidente que o primeiro acto é a razão do segundo e naturalmente o procede (\emph{C}. \emph{Gent}., l II, c. 1) --- Mas o \emph{acto} não se attribue a \textsc{Deus} e ás creaturas num sentido \emph{wnivoco}, mas só num sentido \emph{analogico}; e por isso do acto de \textsc{Deus} devemos afastar toda a imperfeição, que se encontra no acto das creaturas. O acto das creaturas sahe do agente e é recebido no paciente, porque é uma coisa \emph{accidental}, e por isso é formalmente \emph{transeunte}; ao passo que o acto divino não sahe de \textsc{Deus} para as creaturas, porque o acto de \textsc{DEUS} é a sua propria \emph{substancia}, e é simplesmente \emph{inmanente}. --- Todavia o acto divino não é tão \emph{imnmanente}, que se não possa chamar, de algum modo, \emph{transeunte}; porque convem com o acto \emph{éransewnte} das creaturas em que produz um effeito fóra do agente. Por isso, o acto divino chama-se \emph{formalmente} immanente e \emph{eminente} ou \emph{virtualmente} transeunte, emquanto, ficando no agente com que se identifica, produz o que o agente creado não pode produzir senão por um acto, que d'elle passa para o paciente. Emquanto \textsc{Deus} entende e quer, produz o effeito. --- Considerada nas creaturas, a \emph{creação} é uma \emph{relação das creaturas com o Creador}, \emph{como com o principio do proprio ser}. Esta relação é \emph{ndo-mutia}; porque, se as creaturas se referem \emph{realmente} a \textsc{DEus}, \textsc{Deus} só \emph{logicamente} refere-se ás creaturas; pois o \emph{fundamento} da relação, que é \emph{o ser tirado do nada}, só é \emph{real} nas creaturas. Comtudo, não devemos d'ahi concluir que \textsc{Deus} só \emph{logicamente} é Senhor das creaturas; porque, para \textsc{DEus} ser \emph{realmente} Senhor das creaturas, basta que estas \emph{realmente} derivem e dependam de \textsc{DEus}. --- Outras noções, relativas á \emph{creação}, foram expostas na \emph{Cosmologia}, tratando da \emph{origem do mundo}. Ahi consideramos a creação mais no aspecto \emph{passivo}, e aqui mais no aspecto \emph{activo}.}.

\pg{514} 86. A \textbf{creação} é o \textbf{acto} proprio de \textbf{Deus}.

a) A \emph{creacdo}, por ser um acto \emph{absoluto}, isto é, independente de toda a materia, deve ser exclusivamente proprio da causa; que possue uma existencia \emph{absoluta}, isto é, independente de todo o ser; porque o modo de operar é proporcionado ao modo de existir. Ora só \textsc{Deus} possue uma existencia \emph{absoluta}. Logo a \emph{creação} é o acto proprio de \textsc{Deus}.

b) Todo o effeito deve ser proporcionado e semelhante á forma do agente; porque todo o agente opera pela sua forma. Ora a forma, por que \textsc{Deus} opera, é a sua essencia, e esta consiste em que \textsc{Deus} é o proprio \emph{Ser} subsistente, contendo em si toda a plenitude e perfeição da entidade. Logo o effeito proporcionado e semelhante a \textsc{Deus} é o proprio \emph{ser} das coisas, Mas o ser das coisas, pelo facto de abranger \emph{todos} os elementos constitutivos, só pode ser produzido por \emph{creação}. Logo a \emph{creação} é o acto proprio de \textsc{Deus} \footnote{Desenvolvamus e esclareçamos este ponto. A Causa primeira, que communica \textbf{alentras} coisas todo o ser, que possuem, e por isso produz o ente emquanto é \textbf{enigi\&opera} pela \emph{intelligencia} com o concurso da \emph{vontade} e da \emph{potencia operativa}. isso, para entendermos, de algum modo, a acto creador\pgn{515} de \textsc{DEus}, devemos fallar da sua \emph{potencia}, ou \emph{energia operativa}, tendo, no capitulo antecedente, fallado da sua intelligencia e vontade.\par A \emph{potencia}, sendo uma perfeição \emph{absoluta}, deve encontrar-se em \textsc{DEUS} no seu conceito \emph{formal}, \emph{e} em toda a extensão da sua perfeição, excluindo todas as imperfeições e limitações, de que se acha cercada nas creaturas. Ora a \emph{potencia}, no seu conceito \emph{formal}, \emph{é um principio efficiente do effeito}. Este conceito, como se vê, importa simplesmente o \emph{acto}, e exclue toda a imperfeição da \emph{potencialidade}. --- Mas a \emph{potencia} attribue-se só \emph{analogicamente} a \textsc{Deus} e ás creaturas. Por isso, nas creaturas, a \emph{potencia} distingue-se \emph{realmente} tanto da \emph{essencia}, como do \emph{acto}, \emph{é principio efficiente} do \emph{acto} e do \emph{effeito}, e está sempre unida com a potencialidade, porque suppõe a passagem da potencia para o acto; ao passo que, em \textsc{Deus}, no qual a \emph{essencia}, a \emph{potencia} e O \emph{acto} são uma e a mesma coisa (pois \textsc{DEus} é o proprio Ser subsistente), a \emph{potencia operativa}, emquanto ás operações \emph{transeuntes}, não é principio efficiente do \emph{acto}, mas unicamente do \emph{e/feito}, e, d'este modo, encontra-se em \textsc{Deus} em toda a sua perfeição. (\emph{Cf}. \emph{Sum}. \emph{Th}., p. I, q. 25, a. 1 ad 8).\par Emquanto ás relações da \emph{potencia operativa}, ou \emph{activa}, de \textsc{DEUS} com a sua \emph{intelligencia} e \emph{vontade}, dizemos que, na producção e governo das creaturas, a intelligencia \emph{dirige}, á vontade \emph{manda} e a potencia \emph{executa}. Por quanto, \textsc{DEUS}, operando fóra da sua essencia não por necessidade da sua natureza, mas pela liberdade do seu juizo, não produz todos.os effeitos, que a sua potencia é capaz de produzir; mas, como convém a um agente independente, que livremente se determina a produzir um effeito preconcebido pela intelligencia, produz, sem alguma mudança de si mesmo, os effeitos externos, e os produz no modo, na ordem e no tempo, que o eterno juizo da sua vontade livremente determinou. Por isso, todos os effeitos attribuem-se, como á causa efficiente, ao imperio da vontade divina, porque a potencia só produz os effeitos, que a vontade imperou. Com razão diz o Psalmista (\emph{Ps}. 148): u \emph{Ipse dicit et facta sunt}; \emph{Ipse mandavit et creata sunt} » (Cf. S. Th., 2 \emph{Sent}. Dist. I, q. 1, a. 5 ad 11).\par A potencia divina, identificando-se com a \emph{essencia}, é simplesmente \emph{infinita} (\emph{Sum}. \emph{Th}., p. I, q. 25, a. 2). --- E pode chamar-se \emph{infinita} tambem com relação aos \emph{e/feitos}, --- não emquanto a potencia divina produza uma multidão actualmente infinita, ou produza um effeito infinitamente perfeito (o que repugna á razão de effeito) --- mas emquanto não pode exgottar-se na produec\&o dos effeitos, e a efficacia da sua acção nas creaturas é multiplice e illimitada. Diz S. Thomaz: « Numquam (Deus) tot effectus facit, quin plures facere possit, nec umquam ita intense operatur, quin intensius operari possit. Intensio autem in operatione divina non est attendenda secundum quod operatio est in operante, quia sic semper est infinita, cum operatio sit divina essentia; sed attendenda est secundum quod attingit effectum. Sic \textbf{enim} a Deo moventur quaedam efficacius, quaedam minus efficaciter » (\emph{De pot}.,\pgn{516}q. 1, a 2). --- D'onde se segue que \textsc{Deus} é \emph{omnipotente}, emquanto pode fazer todo o que é \emph{absolutamente possivel}. \emph{O absolutamente impossivel} é uma \emph{contradicção}, e \& contradiegüo não pode ser feita. (\emph{Sum}. \emph{Th}., p. I, q. 25, a. 8, --- c. \emph{Gent}., l. II, c. 25, --- \emph{De pot}., q. 1, a. 3).\par Uma potencia \emph{infinita} deve, na sua operação, ser independente de toda e qualquer coisa extranha, e por isso, para produzir o effeito, não exige que alguma coisa d'esse effeito preexista ou se presupponha. O acto, que não presuppõe nada da coisa, que por elle é produzida, diz-se \emph{creacdo}. --- A \emph{creação} é o acto proprio, especificamente proprio de \textsc{Deus}; porque só \textsc{Deus}, por ser a Causa universal e o primeiro Principio das coisas, pode, na sua operação, ser independente de toda e qualquer causa anterior, material ou efficiente. Por isso é absurdo dizer que a creatura possa, por si, tirar do \emph{nada} as coisas; porque é absurdo que a creatura possa ser o primeiro principio das coisas.\par Faz-se a questão se \textsc{DEUS} possa communiocar á creatura a virtude ou força de crear. Os escholasticos respondem que essa virtude ou força não pode ser communieada á creatura, e dão as seguintes razões. --- A creação exige uma potencia \emph{infinita}, que opere pelo \emph{ser}, como tal, e não pelo ser \emph{determinado} a uma certa especie; porque o termo da creacüo é o \emph{ser}, e a causa deve ser proporcionada ao effeito. Ora nenhuma creatura pode receber uma qualidade infinita, porque uma qualidade recebida deve accommodar-se á capacidade do sujeito, que a recebe. --- Além d'isso, o acto da creatura é uma coisa \emph{accidental}, que deve ser necessariamente recebido num sujeito para que possa -produzir o effeito, e por isso a creatura não pode fazer nada senão transformando a materia preexistente: o que é diametralmente opposto ao acto -\emph{creador}. Portanto a creação exige uma força, que se identifique com a essencia do agente, que seja acto puro e não seja distincto da propria operação. (\emph{De pot}., q. 8, \emph{a} 4). --- S. Thomaz accrescenta que a creatura não pode intervir, nem como causa \emph{Instrumental}, no acto \emph{creador} de \textsc{DEus}. De facto, o \emph{instrumento} serve para dispór o sujeito, de modo que este possa receber em si o acto do agente principal; ora, como a creação não presuppõe algum «sujeito, a creatura não pode intervir, nem como \emph{instrumento}, no acto creador de \textsc{Deus}. Além d'isso, o acto de uma causa, embora \emph{instrumental}, deve derivar da potencia d'ella; ora, sendo impossivel que da potencia de uma causa -\emph{finita} derive um acto de uma força infinita, qual é o acto necessario para a creação, é tambem impossivel que a creatura seja causa \emph{instrumental} do "ente \emph{creado} (Cf. \emph{De pot}., q. a. 9, 4; \emph{Sum}. \emph{Th}., p. I, q. 45, a. 5).\par Dissemos que, na producção das creaturas, a intelligencia de \textsc{Deus} \emph{di}- -\emph{rige}, a sua vontade \emph{manda} e a sua potencia \emph{executa}. Para entendermos, de algum modo, o acto da intelligencia divina nessa obra, lembremo-nos da nobilissima theoria do Angelico Doutor. A intelligencia divina, diz S. Thomaz, conhece perfeitamente a essencia divina --- acto purissimo e oceano infinito de toda a perfeição, --- e conhece-a não só emquanto é tal essencia, mas tambem emquanto pode ser imitada e participada por infinitos modos, embora todos infinitamente inadequados. Conhecendo assim a essencia divina, a intelligencia forma innumerayeis conceitos, ou archetypos, de coisas possiveis\pgn{517}, cada um dos quaes imita, mais ou menos, a mesma essencia, e, em conformidade com estes archetypos, \textsc{DEus} produz, numa certa differenca de tempo, tudo o que Elle quer produzir. E 8. Agostinho já tinha ensinado que todas as creaturas existiam, desde a eternidade, na arte e na sapiencia da intelligencia divina, pela qual foram depois creadas: « Sapientia Dei, per quam facta sunt omnia, secundum artem continet omnia, antequam fabricet omnia » (In \emph{Joan}, c. I, n. 17; \emph{cf}. \emph{Sum}. \emph{Th}., p. I, q. 15, a. 4).}.

\pg{515} 87. \textbf{Deus} creou \textbf{todas as} coisas.

a) \emph{Crear} é tirar do \emph{nada}. Ora \textsc{Deus} tirou do \emph{nada} todas as coisas; porque todas as coisas são \emph{contingentes} na \emph{totalidade} da sua substancia, isto é, tanto em relação á \emph{materia} como em relação á \emph{forma}, e todo o ente \emph{contingente} é produzido, em ultima analyse, pelo ente \emph{necessario}. Logo \textsc{Deus} \emph{creou} todas as coisas.

\pg{516} b) Se o acto da Causa primeira, na producção das coisas, devesse presuppór alguma materia n\&o-produzida, esta materia possuiria o ser \emph{por essencia}; porque o ente, que não possue o ser \emph{por participação}, possue-o \emph{por essencia}. Ora é impossivel

\emph{Z} \pg{517}  que a materia possua o ser \emph{por essencia}; porque o ente que possue o ser \emph{por essencia é acto purissimo}, emquanto a materia é \emph{mera potencia}. Logo o acto da Causa primeira, na producção das coisas, não presuppõe alguma materia n\&o-produzida, e por isso é um acto \emph{creador}. Logo \textsc{Deus} \emph{creou} todas as coisas \footnote{Outros argumentos encontram-se em S. Thomaz (C. \emph{Gentes}, 1. II, c. 15, --- \emph{De pot}, q. a. 5). --- Contra esta these fazem-se muitas objecções, algumas das quaes foram expostas e\textasciicircum{}resolvidas na \emph{Cosmologia}. Para a solução de outras, que se referem 4 Causa creadora, expomos os seguintes principios: --- 1º) \textsc{Deus} opera pelas ideas da sua intelligencia, e por isso, assim como não tem propriamente a idea da materia, separada da forma, mas \emph{so} a idea do composto subsistente, assim tambem não crea a materia sem a forma, mas crea a materia formada, isto é, o composto subsistente. Todavia, pode dizer-se que \textsc{Deus} tem uma idea da materia, emquanto a materia imita, de algum modo, a essencia divina (\emph{De pot}., q. 3, a. 1 ad 13). --- 2.9) O acto de \textsc{Deus}, na producção das coisas, passa do agente para o paciente --- emquanto, permanecendo na essencia divina, com a qual se identifica, produz fôra de \textsc{DEus} toda a substancia do paciente, --- mas não emquanto, derivando da essencia divina, é recebido no paciente; porque, como dissemos, o acto divino não se distingue da essencia, e esta produz o effeito sem produzir o acto. --- 8.º) A causa universal que é \textsc{Deus}, produz o ente \emph{emquanto ente} --- não no sentido que produz a razão abstracta do ente communissimo e analogico (que n\&o pode existir na realidade) --- mas no sentido que produz o ente na sua \emph{totalidade}, isto é, segundo \emph{todos} os elementos, de que o ente é constituido e que são a materia e a forma.\par Não nos demoramos neste assumpto, pois que foi sufficientemente tratado na \emph{Cosmologia}. --- Mas não podemos deixar de fazer uma breve critica do systema de A. \emph{Rosmini} acerca da \emph{creacdo}.\par Este escriptor, na sua \emph{T/eosophia}, explica a \emph{creacdo} por um triplice acto da intelligencia divina: \emph{abstracção}, \emph{imaginação} e \emph{synthese}. --- A Pessoa divina do \emph{Pae}, contemplando o \emph{Verbo}, \emph{abstrahe} d'Elle o \emph{ser inicial}, em que vé contido virtualmente todo o ser finito. E' a primeira \emph{creacdo}, que dà origem á primeira creatura, isto á, ao \emph{ser inicial}. Este ser não tem uma existencia \emph{real}, mas somente \emph{logica}, e reside primeiramente na intelligencia de \textsc{Deus} e depois na das creaturas. --- A esta \emph{abstracção} segue-se a \emph{imaginação}, pela qual DEus-PAE concebe varios \emph{termos} reaes e limitados, de que se compõe a realidade do universo. --- Finalmente, realiza-se a \emph{synthese}, pela qual \textsc{Deus} une os dois elementos, que são o \emph{ser inicial} e os diversos \emph{termos} d'elle, isto é, os entes \emph{reaes finitos} (\emph{Teos}., v. 1, ]. 2, sez. 4. c. 6, a. 7).\par O systema de \emph{Rosmini} é falso pelas seguintes razões:\par 1.9) Depois de attribuir só 4 Pessoa do \emph{Pae} a creação, que (como toda a acção de \textsc{Deus} \emph{ad extra}) \emph{é commun}. ás tres divinas Pessoas, \emph{Rosmini} reduz\pgn{518}a oreacüo a um puro \emph{idealismo}. Na verdade, as tres operações, de que ella seria constituida e que são a \emph{abstracção}, a. \emph{imaginação} e a \emph{synthese}, são productos da intelligencia e não dão origem a nenhuma realidade. Porquanto, o proprio \emph{Rosmini} ensina que a primeira creatura --- o \emph{ser inicial} --- não tem existencia \emph{real}, mas só \emph{logica}. Se não tem existencia \emph{real} o producto da \emph{abstracção}, nem a podem ter os productos da \emph{imaginação} e da \emph{synthese}. D'este modo, os entes creados não passam de entes \emph{ideaes}.\par 2.9) \emph{Rosmini} admitte como \emph{increado} o ser das creaturas. De facto, o ser das creaturas, na opinião d'esse escriptor, é constituido pelo \emph{ser inicial}, e este \emph{ser inicial} é unico, simplicissimo e eterno, porque é o proprio \emph{ser} do \emph{Verbo}, abstrahido somente da subsistencia. Ora esta opinião, além de conter o \emph{pantheismo}, destroe o conceito da \emph{creação}; porque a creação, como dissemos, tem por termo o proprio \emph{ser}, pois produz a \emph{totalidade} da substancia, Assim ensina S. Thomaz: « Oreatio terminatur ad \emph{esse}, tamquam ad proprium effectum » (\emph{De pot}., q. 8, a. 16 ad 21) E noutro logar: « Cum Deus sit \emph{ipsun esse} per suam essentiam, oportet quod \emph{esse creatum} sit proprius effectus ejus, sicut ignire est proprius effectus ignis » (\emph{Sum}. \emph{Th}., P. Lq. 8, a. 1). --- Goherente com esse principio, \emph{Rosmini} diz que --- o ser das creaturas não é \emph{produzido}, mas é \emph{posto}, --- o \emph{ser das creaturas} não se distingue do \emph{ser divino}, senão por causa dos \emph{termos}, --- só \textsc{Deus} pode crear, porque o \emph{ser inicial}, que dá á creatura a forma de ente, é \emph{alguma coisa} de \textsc{Deus}, e só \textsc{DEus} pode dispór do que lhe pertence (\emph{Teos}., v. 1, p. 969).}.

\pg{518} 88. \textbf{Deus} foi \textbf{livre na creação} do mundo. --- \textsc{Deus} foi \emph{livre} na creacüo do mundo, se a-creação não era exigida --- nem pela essencia divina, --- nem pela essencia do mundo. Ora a creação do mundo não era exigida --- nem pela essencia divina, porque esta é um oceano infinito de perfeições, e, para ser felieissima, não precisa de algum bem creado, --- nem pela essencia do mundo, porque o mundo é contingente, e todo o ente contingente é por si indifferente para a existencia e para a não-existencia. Logo \textsc{Deus} foi \emph{livre} na creação do mundo \footnote{Não insistimos neste ponto, que foi sufficientemente desenvolvido, tratando da \emph{vontade} de \textsc{Deus}. \textsc{Deus} é causa do mundo, --- não \emph{natural}, isto é, que não opera por necessidade da sua natureza, --- mas \emph{livre}, isto é, que opera pelo livre juizo da sua vontade. Por isso, não só não foi necessario que \textsc{DEus} creasse o mundo \emph{desde a eternidade}, mas nem foi necessario que o creasse \emph{no tempo}. Diz S. Thomaz: « Voluntas Dei est causa rerum. Sic ergo aliqua necesse est esse, sicut necesse est Deum velle illa: cum necessitas effeotus ex necessitate causae dependeat. Ostensum est autem quod, absolute loquendo, non est necesse Deum velle aliquid, nisi seipsum. Non est ergo necessarium Deum velle quod mundus fuerit semper. Sed eatenus mundus est, quatenus Deus vult illum esse, cum esse mundi ex voluntate Dei dependeat, sicut ex sua causa. Non est \emph{igitur} necessarium mundum Semper esse» (\emph{Sum}. \emph{Th}., p. 1, q. 46, a 1). --- Lembramos que, quando se diz que Deus creou o mundo \emph{no tempo}, --- não entendemos que o tempo tenha preexistido ao mundo, ou que o tempo seja a medida de toda a duração, mesmo da duração divina (o que é absurdo), --- mas entendemos que Deus, immovel na sua eternidade, que abrange e excede toda a duração\pgn{519}temporanea, coreou o mundo \emph{com o tempo}, emquanto, creado o mundo, começou o \emph{movimento}, de que o tempo é medida. O tempo, que se pensa ter existido antes do mundo, é \emph{Zmaginario}, não \emph{real} (\emph{Of}. \emph{De pot}., q. 9, a. 17).}.

\pg{519} 89. \textbf{Deus creou} o mundo para a sua gloria. --- \textsc{Deus}, sendo infinitamente intelligente, devia propor-se um fim na creação do mundo. Este fim não pode ser senão a communicação e por isso a manifestação das perfeicoes divinas, --- communicação e manifestação, de que resulta a \emph{gloria} de \textsc{Deus}. Logo \textsc{Deus} ereou o mundo para a sua \emph{gloria} \footnote{A \emph{gloria} é o \emph{conhecimento} claro da \emph{bondade} de um ser, a qual é digna de louvor. Por isso define-se: \emph{clara notitia cum laude}. D'onde se segue que a \emph{gloria} é maior ou menor, conforme maior ou menor é a \emph{bondade}, de que um ser é dotado, e conforme maior ou menor é o \emph{conhecimento}, que se tem d'essa bondade. A gloria das creaturas é imperfeita, porque a sua bondade é finita e participada. Só a gloria de \textsc{Deus} é infinitamente perfeita, porque só a sua bondade é essencial e infinita. --- Mas, como a bondade divina só por \textsc{DEus} pode ser conhecida adequadamente e por isso louvada como merece, só a gloria, que \textsc{Deus} recebe de si mesmo, é digna de \textsc{DEus}. A gloria, que a creatura pode e deve dar a \textsc{Deus}, é sempre infinitamente inferior ao merito; porque infinitamente deficiente é o conhecimento, que a creatura tem da bondade divina. --- Ainda que infinitamente inferiores ao merito, os nossos louvores dão ao bom \textsc{Deus} uma satisfação immensa e procuram á nossa alma uma felicidade eterna.\par Esta proposição --- que \textsc{Deus} ereou o mundo para a \emph{sua gloria} --- não é contraria á outra proposição --- que \textsc{DEus} ereou o mundo para communicar a sua bondade. Communicando a sua bondade, \textsc{Deus} entende a sua gloria. Como tambem, dando \emph{gloria} a \textsc{DEus}, as creaturas participam da felicidade divina. \textsc{DEus} não quiz separar do nosso bem a sua gloria; o nosso bem é a sua gloria, e a sua gloria é o nosso bem. Felizes as almas, que reconhecem os beneficios de \textsc{Deus} e Lhe tributam a homenagem da gratidão, fidelidade e da obediencia!}.
